\section{Peano Arithmetics}

\titledbox{Remember}{
A unary natural number is a term generated by the following grammar rule:

\[
u ::= 0 \mid s(u)
\]

where $0$ denotes zero and $s(u)$ denotes the successor of $u$. Therefore, the set of unary natural numbers is:
\[
\set{0, s(0), s(s(0)), s(s(s(0))), \cdots}
\]
}

\begin{excursion}[Peano Axioms]{}
The following set of axioms are called the Peano axioms. They can be used to \textit{reason} about natural numbers:
\begin{itemize}
    \item \textbf{P1.} $0$ is a unary natural number.
    \item \textbf{P2.} Every unary natural number has a successor that is a unary natural number and that is different from it.
    \item \textbf{P3.} $0$ is not a successor of any unary natural number.
    \item \textbf{P4.} Different unary natural numbers have different successors.
    \item \textbf{P5.} Induction Axiom. Every unary natural number possesses a property $P$ if
    \begin{itemize}
        \item $0$ has the property $P$, and
        \item the successor of every 
        unary natural number that has the property $P$ also possesses property $P$
    \end{itemize}
\end{itemize}
\end{excursion}

\begin{excursion}[Unary Operations]{}
\definition[Addition]{
We define the addition operation inductively:
\begin{itemize}
    \item Adding $0$ to a number does not change it, expressed as: $0 + X = X$.
    \item Adding the successor of $X$ to $Y$ yields the successor of $X + Y$, expressed as: $s(X) + Y = s(X + Y)$.
\end{itemize}
}{}

Similarly, multiplication and exponentiation are defined as follows:
\definition[Multiplication]{
\begin{itemize}
    \item $n \cdot 0 = 0$, and
    \item $n \cdot s(m) = n + (n \cdot m)$
\end{itemize}
}{}

\definition[Exponentiation]{
\begin{itemize}
    \item $n^0 = s(0)$, and
    \item $n^{s(m)} = n \cdot n^m$
\end{itemize}
}{}
\end{excursion}

In Prolog, we can represent unary natural numbers with a predicate \hltext{nat/1} and a functor \hltext{s/1}:

\begin{codeblock}
nat(0).             % 0 is a natural number
nat(s(X)) :- nat(X).  % the successor of a natural number is a natural number
\end{codeblock}

This predicate will return true for all the following:

\begin{codeblock}
nat(0), nat(s(0)), nat(s(s(0))), nat(s(s(s(0)))), ... 
\end{codeblock}

\questionbox{
Can we write unary addition?
}

Let's think about it. This is a \textbf{function} that returns some value given two arguments, so we have to represent that as a $3$-ary predicate.

Suppose we want to do that with a predicate \hltext{add/3}

\begin{codeblock}
add(X,Y,Z)        %true if Z is the result of adding X and Y
\end{codeblock}

Unary addition is recursive. The base case is:

\[
0 + X = X
\]

We simply write:
\begin{codeblock}
add(0,X,X).
\end{codeblock}

The recursive case:

\[
s(X) + Y = s(X + Y)
\]

In Prolog, we declare the third argument (the output) as something of the structure \inlinecode{s(Z)} and in the body we compute that $Z$ as the addition of $X$ and $Y$.

\begin{codeblock}
add(s(X), Y, s(Z)) :- add(X, Y, Z).
\end{codeblock}

Small test ($3 + 2$):

\begin{codeblock}
?- add(s(s(s(0))), s(s(0)), X).
X = s(s(s(s(s(0)))))
\end{codeblock}

\questionbox{
Can we do subtraction?
}

\titledbox{Beauty of Prolog}{
We do not need a separate subtraction predicate! We arranged the arguments in our head to mean addition. We can query the \inlinecode{add/3} predicate itself to have answers for subtractions. Because \inlinecode{add(X, Y, Z)} defines the mathematical relation, subtraction comes for free without writing any new code.
}

\begin{example}{}
We want the result of $3-1=2$. We query our knowledge base to look for the appropriate $X$ such that $X+1=3$:

\begin{codeblock}
?- add(X, s(0), s(s(s(0)))).
X = s(s(0))
\end{codeblock}
\end{example}

Of course, we can define it recursively if we want. Subtraction can be seen as:

\begin{itemize}
    \item $X - 0 = X$ \hfill{Base Case}
    \item $s(X) - s(Y) = X-Y$ \hfill{Recursion}
\end{itemize}

\begin{codeblock}
% Base case: subtracting 0 from any number X leaves X
sub(X, 0, X).

% Recursive step: (X + 1) - (Y + 1) is the same as X - Y
sub(s(X), s(Y), Z) :-
    sub(X, Y, Z).
\end{codeblock}

Small test ($3-1=2$):
\begin{codeblock}
?- sub(s(s(s(0))), s(0), Res).
Res = s(s(0))
\end{codeblock}

And we can define it in terms of addition:

\begin{codeblock}
% sub(A, B, Res) is true if A - B = Res, which means Res + B = A
sub(A, B, Res) :- add(Res, B, A).
\end{codeblock}


Similarly, we can represent unary multiplication:

\[
X \cdot 0 = 0 \qquad X \cdot s(Y) = X + (X \cdot Y)
\]

\titledbox{Sanity Check}{
\[
4 \cdot 3 = 12 \qquad 4 \cdot s(2) = 4 + (4 \cdot 2) = 4 + 8 = 12
\]
}

We write a predicate \inlinecode{mult/3}:

\begin{codeblock}
% Base case
mult(_, 0, 0).

% Recursive step
mult(X, s(Y), Z) :-
    mult(X, Y, W),
    add(X, W, Z).
\end{codeblock}

To see how this works, consider the query computing $2 \times 2 = 4$:
\begin{codeblock}
?- mult(s(s(0)), s(s(0)), Ans).
\end{codeblock}

The execution proceeds by recursively decomposing the second argument down to the base case, and then accumulating the sum on the way back up:

\begin{enumerate}[label=\textbf{Step \arabic*:}, leftmargin=*]
    \item The first call \texttt{mult(s(s(0)), s(s(0)), Ans)} matches the recursive rule with:
    \[
    X = s(s(0)), \quad Y_1 = s(0)
    \]
    To compute the answer, Prolog creates two subgoals:
    \begin{enumerate}[label=(\alph*)]
        \item Compute the sub-product: \texttt{mult(s(s(0)), s(0), $W_1$)}.
        \item Then add: \texttt{add(s(s(0)), $W_1$, Ans)}.
    \end{enumerate}

    \item The second call is the first subgoal \texttt{mult(s(s(0)), s(0), $W_1$)}. This matches the recursive rule again with:
    \[
    X = s(s(0)), \quad Y_2 = 0
    \]
    It generates two further subgoals:
    \begin{enumerate}[label=(\alph*)]
        \item Compute: \texttt{mult(s(s(0)), 0, $W_2$)}.
        \item Then add: \texttt{add(s(s(0)), $W_2$, $W_1$)}.
    \end{enumerate}

    \item Prolog evaluates \texttt{mult(s(s(0)), 0, $W_2$)}. Since the second argument is $0$, it directly matches the \textbf{base case} \inlinecode{mult(_, 0, 0)}, giving:
    \[
    W_2 = 0
    \]

    \item \textbf{Unwinding the Additions:}
    Now the engine resumes and resolves the pending additions back up the call stack:
    \begin{itemize}
        \item \textbf{Resolving $W_1$:} \texttt{add(s(s(0)), 0, $W_1$)} yields:
        \[
        W_1 = s(s(0)) \quad (\text{i.e., } 2 + 0 = 2)
        \]
        \item \textbf{Resolving \texttt{Ans}:} \texttt{add(s(s(0)), s(s(0)), Ans)} yields:
        \[
        \text{Ans} = s(s(s(s(0)))) \quad (\text{i.e., } 2 + 2 = 4)
        \]
    \end{itemize}
\end{enumerate}

All goals succeed, and Prolog returns:
\begin{codeblock}
Ans = s(s(s(s(0))))
\end{codeblock}

\commandnote{
Division comes for free too
}

\begin{example}[Division using Multiplication]{}
Prolog finds $4/2=2$ without writing a division predicate.
\begin{codeblock}
?- mult(s(s(0)), X, s(s(s(s(0))))).
X = s(s(0))
\end{codeblock}
    
\end{example}

Similarly, we can represent unary exponentiation:
\[
X^0 = s(0) \qquad X^{s(Y)} = X \cdot X^Y
\]

\titledbox{Sanity Check}{
\[
3^3 = 3^{s(2)} = 3 \cdot 3^2 = 3 \cdot 9 = 27
\]
}

We write a predicate \inlinecode{uexp/3}:

\begin{codeblock}
% Base case
uexp(_, 0, s(0)).

% Recursive step
uexp(X, s(Y), Z) :-
    uexp(X, Y, W),
    mult(X, W, Z).
\end{codeblock}

Small test ($2^2 = 4$):
\begin{codeblock}
?- uexp(s(s(0)), s(s(0)), Ans).
Ans = s(s(s(s(0))))
\end{codeblock}